\scanpage{8}
% source lines 1-4
On s'intéresse aux sous-groupes \inserted{ouverts} de $C(K)$ \deleted{\unsure{contenant la partie connexe}}, i.e. aux sous-groupes \inserted{ouverts} de $\unsure{J|\mathscr{J}}(K)$ contenant $K^*$. Un tel sous-groupe $M$ contient \unsure{aussi} $J_\infty^{\mathrm{conn}}\deleted{\unsure{\subset}}=J^{\mathrm{conn}}$, \deleted{et} un sous-groupe \deleted{$\unsure{H_{\underline m}V_\infty}$} $V_{\underline m}$ [où

% source lines 5-8
\diagram{\begin{tikzcd}[ampersand replacement=\&,column sep=0.5em,row sep=small]
\bigl(\quad V_{\underline m}=\prod_{\substack{\unsure{\mathfrak p\nmid\underline m|\mathfrak p\mid\underline m}\\\mathfrak p\text{ fini}}}\&V_{\underline m,\mathfrak p}\arrow[dr]\&\subset\deleted{\unsure{C_a}}\subset{\unsure{D|\mathscr{D}}}_{\mathrm{abs}}\\
\&\&\begin{gathered}\underline m\text{-unités;}\\\text{à distance infinie, cela est }J_\infty^{\mathrm{conn}};\\\text{si }\mathfrak p\nmid\underline m,\text{ cela est }V_{\mathfrak p}\end{gathered}
\end{tikzcd}}
\[
\begin{aligned}
\bigl(\quad&J_{\underline m}=J\cap\prod_{\mathfrak p\nmid\underline m}K_{\mathfrak p}\\
\bigl(\quad&D_{\underline m}=J_{\underline m}/V_{\underline m}\cap J_{\underline m}\simeq\coprod_{\mathfrak p\nmid\underline m}\deleted{\unsure{V}}W_{\mathfrak p}=\mathbb Z^{\pi(\underline m)}
\end{aligned}
\]
($\underline m$ est un diviseur absolu contenant les places à l'$\infty$)
\margin{$\pi(\underline m)$ ens des diviseurs premiers $\nmid\underline m$}
(diviseurs premiers à $\underline m$)
]

% source lines 9-11: cancelled block
\deleted{et un sous-groupe $\displaystyle\prod_{\substack{\mathfrak p\text{ fini}\\\mathfrak p\mid\underline m}}A_{\mathfrak p}\unsure{\simeq|=}A'$ où pour tout $\mathfrak p$ fini, $\mathfrak p\mid\underline m$, $A_{\mathfrak p}$ est un voisinage de $1$ dans $K_{\mathfrak p}$.}

% source lines 12-18
Donc $M\supset\deleted{\unsure{H_{\underline m}=(V_{\underline m}\times\prod_{\mathfrak p\mid\underline m,\ \mathfrak p\text{ fini}}A_{\mathfrak p}\times J_\infty^{\mathrm{conn}})}}V_{\underline m}K^*$, et réciproquement, si \inserted{\unsure{$M$ contient un $H_{\underline m}$}}, ....

Je dis qu'on a $\deleted{\unsure{H_{\underline m}}}K^*V_{\underline m}+J_{\underline m}=J$, ce qui \unsure{résulte|résultera} du fait que, \deleted{$\unsure{V_{\underline m}\in A'\times J_\infty^{\mathrm{conn}}}$} \inserted{$V_{\underline m}\cap\displaystyle\prod_{\mathfrak p\mid\underline m}K_{\mathfrak p}$} est un voisinage de $1$ dans $\deleted{\unsure{J}}\displaystyle\prod_{\mathfrak p\mid\underline m}K_{\mathfrak p}$, il existe pour tout $\vec x_{\unsure{\mathfrak p|\underline m}}\in\displaystyle\prod_{\mathfrak p\mid\underline m}K_{\mathfrak p}$ un $a\in K^*$ tel que $a_{\mathfrak p}^{-1}\vec x\in V_{\underline m}\deleted{\unsure{\times V_\infty}}$; Ainsi,

% source lines 19-21
On aura
\[
J/K^*V_{\underline m}\simeq J_{\underline m}/K^*V_{\underline m}\cap J_{\underline m}
\]
\deleted{\unsure{on}} et \deleted{\unsure{comme}} d'autre part, comme $H_{\underline m}\supset V_{\underline m}\cap J_{\underline m}$, \deleted{\unsure{on}} on a
\[
J/H_{\underline m}\simeq D_{\underline m}(K)/P_{\underline m}(K)
\]

% source lines 22-24 and lower margin
où $P_{\underline m}(K)=$ \unsure{ens|\deleted{les}} des diviseurs \deleted{\unsure{principaux}} \inserted{premiers à $\underline m$} de la forme $\operatorname{div}(\vec y)$, i.e. $\vec y=a\cdot\vec x\deleted{\unsure{\vec x_{\underline m}\times\vec x_\infty}}$ avec

\deleted{$a\in K^*,\quad\vec x'_{\underline m}\times\vec x''_{\underline m}\times\vec x_\infty\in(J_{\underline m}\cap V_{\underline m})\times\unsure{\vec x'_\infty}$}

$a\in K^*$, $\vec x\in V_{\underline m}$, $\deleted{\unsure{\operatorname{div}(\vec y)\in D_{\underline m}}}\quad y_{\mathfrak p}=1\text{ si }\mathfrak p\mid\underline m$.
% slanted
\margin{diviseurs \unsure{non} absolus}
% archive sheet number: 7
